\chapter{Installation}
\label{ch:installation}

OmniLaTeX can be set up using several methods, each suited to a different
workflow.  This chapter covers the recommended Docker approach, the Nix
alternative, a manual TeX~Live installation, and a one-click DevContainer setup.

\section{Prerequisites}

Before choosing an installation method, review the requirements for each
approach:

\begin{description}
    \item[Docker (recommended)] Docker Engine 20.10 or later.  The OmniLaTeX
        image bundles a complete TeX~Live installation, Python tooling,
        fonts, and all external dependencies.  No separate \TeX{} installation
        is required on the host.
    \item[Nix (alternative)] Nix with Flakes enabled.  The \texttt{flake.nix}
        declares all \TeX{} packages explicitly, producing a fully
        reproducible development shell.
    \item[TeX~Live 2026 (manual)] A full TeX~Live installation including
        \hologo{LuaLaTeX}, Biber, \texttt{bib2gls}, Inkscape, Gnuplot, and
        Python~3 with Pygments.  Most users should prefer Docker or Nix.
\end{description}

All methods require Git to clone the repository:

\begin{minted}{bash}
git clone https://github.com/WyattAu/OmniLaTeX-template.git
cd OmniLaTeX-template
\end{minted}

\section{Docker Method}

The Docker image is the recommended way to use OmniLaTeX.  It provides a
self-contained build environment with TeX~Live, system fonts (Libertinus,
Monaspace, Atkinson Hyperlegible Next), and all auxiliary tools pre-installed.

\subsection{Pulling the Image}

\begin{minted}{bash}
docker pull ghcr.io/wyattau/omnilatex-docker:latest
\end{minted}

The image is published to the GitHub Container Registry (ghcr.io).  The
\texttt{latest} tag tracks the most recent TeX~Live release.

\subsection{Verifying the Installation}

Confirm the image works by invoking \texttt{latexmk} through the container
entrypoint:

\begin{minted}{bash}
docker run --rm -v "$(pwd):/workspace" \
    ghcr.io/wyattau/omnilatex-docker:latest --version
\end{minted}

The entrypoint script at \texttt{/usr/local/bin/entrypoint.sh} forwards
arguments to \texttt{latexmk} by default.  This means you can compile any
document by passing its filename directly.

\subsection{Running Your First Build}

Mount the repository into the container and run the full build:

\begin{minted}{bash}
docker run --rm -v "$(pwd):/workspace" \
    ghcr.io/wyattau/omnilatex-docker:latest \
    python3 build.py build-all --mode dev
\end{minted}

This compiles the root document and all example projects.  The resulting PDFs
are written to the \texttt{build/} directory on the host.

\subsection{Bypassing the Entrypoint}

To execute arbitrary commands inside the container (e.g.\ running tests or
opening a shell), override the entrypoint:

\begin{minted}{bash}
docker run --rm -it --entrypoint bash \
    -v "$(pwd):/workspace" \
    ghcr.io/wyattau/omnilatex-docker:latest
\end{minted}

Inside the shell, all TeX~Live binaries, Python packages, and fonts are on
\texttt{PATH}.

\section{Nix Method}

The \texttt{flake.nix} at the repository root declares the complete OmniLaTeX
development environment as a Nix flake.  This method guarantees reproducibility:
every collaborator gets identical package versions.

\subsection{Entering the Dev Shell}

\begin{minted}{bash}
nix develop .#
\end{minted}

This drops you into a shell with TeX~Live (via \texttt{texlive.combine}),
Python~3 with Pygments and Rich, Inkscape, Gnuplot, and Lean~4 available.

The TeX~Live scheme is \texttt{scheme-medium} plus explicit package
declarations extracted from every \texttt{\textbackslash usepackage} and
\texttt{\textbackslash RequirePackage} call in \texttt{lib/*.sty} and
\texttt{omnilatex.cls}.  This avoids the $\sim$4\,GB \texttt{scheme-full}
while covering all OmniLaTeX dependencies.

\subsection{Building from the Nix Shell}

\begin{minted}{bash}
python3 build.py build-all
\end{minted}

For Lean~4 formal verification targets:

\begin{minted}{bash}
lake build
\end{minted}

\subsection{Building Individual Examples with Nix}

The flake exposes per-example packages.  To build a single example as a Nix
derivation:

\begin{minted}{bash}
nix build .#examples-article
\end{minted}

The resulting PDF is placed at \texttt{result/omnilatex-article.pdf}.

\section{TeX Live Method}

After OmniLaTeX is published on CTAN, the class file and all bundled style
files can be installed via the TeX~Live package manager:

\begin{minted}{bash}
tlmgr install omnilatex
\end{minted}

This installs \texttt{omnilatex.cls} and the \texttt{lib/} style library into
the local TeX~Live tree.  You will still need to install the external tools
listed in Section~\ref{sec:prerequisites} separately.

For development (using the class from the repository without a CTAN release),
set the \texttt{TEXINPUTS} environment variable:

\begin{minted}{bash}
export TEXINPUTS="/path/to/OmniLaTeX-template:"
\end{minted}

\section{DevContainer}

VS Code users can open the repository directly as a DevContainer.  The
configuration in \texttt{.devcontainer/devcontainer.json} launches the
OmniLaTeX Docker image with recommended extensions pre-installed:

\begin{itemize}
    \item \texttt{ms-python.python} -- Python language support.
    \item \texttt{james-yu.latex-workshop} -- \LaTeX{} compilation recipes
          configured for dev and prod modes.
    \item \texttt{streetsidesoftware.code-spell-checker} -- Spell checking.
\end{itemize}

To start:

\begin{enumerate}
    \item Open the repository in VS Code.
    \item When prompted, select \fbox{\small Reopen in Container}.
    \item VS Code builds the container and installs extensions automatically.
\end{enumerate}

The \texttt{postCreateCommand} runs \texttt{python3 build.py preflight} to
verify the environment on first launch.

\section{Verification}

After installation, verify that everything works by running the full test
suite:

\begin{minted}{bash}
python3 build.py test
\end{minted}

To check all example documents compile successfully:

\begin{minted}{bash}
python3 build.py build-examples --mode dev
\end{minted}

A successful run produces PDFs for all 43 examples in the \texttt{build/}
directory.  The build orchestrator checks for PDF existence (not exit codes) to
determine success, since some \LaTeX{} warnings do not constitute failures.

For per-example timing information:

\begin{minted}{bash}
python3 build.py build-all --timings
\end{minted}

\section{Troubleshooting}
\label{sec:installation-troubleshooting}

Table~\ref{tab:installation-issues} lists the most common installation
problems and their solutions.

\begin{table}[htbp]
    \centering
    \caption{Common installation issues and resolutions}
    \label{tab:installation-issues}
    \begin{tabular}{lp{7.5cm}}
        \toprule
        \textbf{Symptom} & \textbf{Resolution} \\
        \midrule
        Font not found (Libertinus, Monaspace) &
            Use the Docker image, which bundles all custom fonts.  For Nix,
            install fonts system-wide and regenerate the fontconfig cache with
            \texttt{fc-cache -f -v}. \\[6pt]
        Permission denied on \texttt{build/} &
            Docker creates files as UID/GID 1000 by default.  If your host
            UID differs, pass \texttt{-e USER\_ID=\$(id -u)} or adjust volume
            permissions. \\[6pt]
        Out of memory during compilation &
            TeX~Live with \ctanpackage{tikz} and \ctanpackage{pgfplots} can
            require $>$2\,GB of RAM.  Increase Docker's memory limit in
            Desktop settings, or close other applications. \\[6pt]
        \texttt{pygmentize} not found &
            The \ctanpackage{minted} package requires Python's Pygments
            library.  Install it with \texttt{pip install pygments} or use
            the Docker image where it is pre-installed. \\[6pt]
        \texttt{bib2gls} not found &
            Required by \ctanpackage{glossaries-extra}.  Install via
            \texttt{tlmgr install bib2gls} or use Docker/Nix. \\[6pt]
        \texttt{luaotfload} font cache error &
            Run \texttt{luaotfload-tool --update} to regenerate the font name
            database.  This is done automatically in the Docker image. \\
        \bottomrule
    \end{tabular}
\end{table}
